\scanpage{17}
\uline{Si} $n+\rho$ \uline{est impair}, alors $b_n=\dim\underbrace{H^{\unsure{n+\rho|n}}(\unsure{\overline{\bar X,\bar E}|\bar X,\bar E})}_{\text{poids }n+\rho}$ est pair, donc $\deleted{\unsure{b_n|b}}(-1)^{b_n}=1$. De plus, un raisonnement du cours de Serre utilisant une forme alternée non dégénérée invariante montre que dans ce cas $\operatorname{signe}\operatorname{dét}f_{H^n(\unsure{X|\bar X},\bar E)}$ est \unsure{également|égal à} $+1$. Donc on trouve alors $\varepsilon'(E)=1$.

Si $n+\rho$ est pair, on ne peut déterminer le signe de façon absolue. On peut \unsure{seult|seul\textsuperscript{t}} \unsure{obs.|obs} \inserted{(indép. de la parité)} \deleted{que} si $\varepsilon_1$ $(\varepsilon_{-1})$ est la \deleted{\unsure{par dimension|dimension|par dimension de}} \inserted{multiplicité} de \deleted{\unsure{l'esp.|l'espace}} la valeur propre $1$ $(-1)$ de $f$ sur $\unsure{H^n|H}(\bar X,\bar F)$, on a $(-1)^{b_n}=(-1)^{\varepsilon_1+\varepsilon_{-1}},\quad\varepsilon(E)=(-1)^{\varepsilon_{-1}}$, donc
\[
\boxed{\varepsilon'(E)=(-1)^{\varepsilon_1}}.
\]

Lorsque l'on n'est pas dans le cas favorable $X$ lisse, $E$ \unsure{ct tordu|ct tordus|de torsion}, $\deleted{\unsure{E|F}}\check E$ \uline{globalement} isomorphe à $E(\rho)$, alors $D(E)$ ne s'exprime plus de façon aussi simple en fonction de $E$, et l'équation fonctionnelle (1) ne s'exprime plus \unsure{néc.|nic.} en termes de la seule fonction $L_E$ (mais elle relie \uline{deux} fonctions $L_{DE},L_E$]
% archive sheet number: 16
